% \usepackage[T1]{fontenc} % font encoding
% \usepackage{fontspec}
\usepackage{fontspec}


\makeatletter%
\@ifclassloaded{beamer}{%
  \usepackage{lmodern} %%% modern beamer style
  %%beamer style
  \usetheme{metropolis}
  \usefonttheme{professionalfonts}
  \setbeamercolor{background canvas}{bg=background}
  \setbeamercolor{frametitle}{bg=snow,fg=black}
  \setbeamercolor{title separator}{fg=snow}
  \setbeamercolor{alerted text}{fg=accent}

  \let\oldtitle\title
  \renewcommand{\title}[1]{\oldtitle{CMPUT391\\#1}}
  \author{Instructor: Denilson Barbosa}
  \date{\small University of Alberta}
  \institute{\today\vskip4em Slides by D. Barbosa, with suggested corrections and improvements by (in alphabetical order) C. Bins, D. Caminhas, K. Guhzva, Q. Lautischer, E. Macdonald, M. A. Nascimento, K. Newbury, M. Strobl, D. Sunderman, K. Wang, and K. Wong.}
  
  \titlegraphic{\includegraphics[width=1.5cm]{../images/by-sa.png}~
  {\tiny This work is licensed under a Creative Commons Attribution 4.0 International License}}

  %%%% un-framed blocks
  \setbeamertemplate{blocks}[rounded][shadow=false]
}{}

\@ifpackageloaded{xcolor}{}{
  \usepackage{xcolor}
}
\makeatother

\setsansfont[BoldFont={Open Sans SemiBold}]{Open Sans}[Scale=0.9]
\setmonofont{Cousine}[Scale=0.875]
% \setmathtt{Cousine}[Scale=0.875]

%%%% highlighting text 
\usepackage{xspace}
\def\highlight#1{\colorbox{accent}{\textcolor{white}{#1}}\xspace}

\def\blue#1{\textcolor{blue}{#1}}

%
%
%
\usepackage{parskip}

\usepackage{subcaption} % needed for subfigures

\usepackage{wrapfig} % for figures "to the side of the text"

\usepackage{anyfontsize} % for egregious warnings

%
% Fonts
%
% \usepackage{amsmath}
% \usepackage{mathspec}
% \setmathtt[Scale=0.875]{Cousine}

%
% various packages for tables and colored tables
%
\usepackage{xcolor,colortbl}
\usepackage{tcolorbox}

\definecolor{accent}{RGB}{225,68,85}
\definecolor{highlight}{RGB}{170,68,153}
\definecolor{background}{rgb}{253,252,252}

\definecolor{sqlColor}{RGB}{119,119,17}

\definecolor{Cfunction}{RGB}{51,51,153}

\definecolor{stringColor}{RGB}{51,51,153}
\definecolor{datatypeColor}{RGB}{170,68,153}
\definecolor{functionColor}{RGB}{51,134,116}

\definecolor{exampleColor}{RGB}{51,34,136}
\definecolor{commentColor}{RGB}{112,68,225}

\definecolor{fern}{RGB}{80,118,66}

\definecolor{snow}{RGB}{240,240,230}
\definecolor{Gray}{gray}{0.85}

\definecolor{tupleBoxColor}{gray}{0.85}

\definecolor{Maroon}{RGB}{139,0,0}

%
% coloring of SQL code
%
\def\ColorSymbol#1{\textcolor{functionColor}{#1}}

%
% hihglight box
%
\usepackage{xspace}
\def\highlight#1{\colorbox{accent}{\textcolor{white}{#1}}\xspace}

%
% Tables
%
\usepackage{multirow}
\usepackage{multicol}

%
% CODE LISTINGS
%
\usepackage{listings}

%
% catch-all style meant to provide a clean and "empty" style
% from which all others build upon
%
\lstdefinestyle{cmput391}{
  style=,
  frame=no,
  tabsize=2,
  numbers=none,
  showstringspaces=false,
  numberstyle=\footnotesize,
  basicstyle=\ttfamily\footnotesize\color{black},
  numbersep=4pt,
  keywords=[1]{},
  keywords=[2]{},
  keywords=[3]{sh_lock,xl_lock,unlock,read,write,commit,abort},
  keywordstyle=[1]\ttfamily\bfseries\footnotesize,
  keywordstyle=[2]\ttfamily\footnotesize,
  keywordstyle=[3]\ttfamily\footnotesize\color{Maroon},
  emph=[1]{},
  emph=[2]{},
  emph=[3]{},
  emphstyle=[1]\ttfamily\footnotesize,
  emphstyle=[2]\ttfamily\footnotesize,
  emphstyle=[3]\ttfamily\footnotesize,
  commentstyle=\ttfamily\itshape\color{commentColor},
  stringstyle=\ttfamily\color{stringColor},
  moredelim=**[l][\itshape\color{commentColor}]{--},
  moredelim=**[is][\color{highlight}]{-|}{|-},
  moredelim=[is][\bfseries\color{highlight}]{-:}{:-},
  morestring=[s][\color{stringColor}]{"}{"},
  morecomment=[l]{--}
}

\lstdefinelanguage{DTD}{
  style=cmput391,
  morekeywords=[1]{DOCTYPE,ELEMENT,ATTLIST},
  keywordstyle=[1]\color{accent},
  literate = {\#PCDATA}{{\textcolor{Cfunction}{\#PCDATA}}}7
             {\#IMPLIED}{{\textcolor{Cfunction}{\#IMPLIED}}}8
}

\lstdefinelanguage{XPath}{
  style=cmput391,
  morekeywords={xs,integer,current,date,distinct,values,avg,text,id,count,position,and,or,not,boolean,number,sum,floor,ceiling,round,doc,eq,ne,gt,lt,ge,le},
  keywordstyle=\color{red},
  literate = {last(}{{\textcolor{functionColor}{last}(}}5 {@<}{{<}}1  
}

\lstdefinestyle{XQuery}{
  style=cmput391,
  keywords=[1]{if,then,else,and,or,some,every,satisfies,for,in,let,where,group,order,by,return,ascending},
  keywords=[2]{gt,lt,eq,ne,ge,le,not,or},
  keywords=[3]{text,sum,min,max,avg,len,position,number,boolean,floor,ceiling,round,%
               id,doc,count,current-date,distinct-values,xs:integer,subsequence},
  keywordstyle=[1]\bfseries\color{sqlColor},
  keywordstyle=[2]\color{accent},
  keywordstyle=[3]\color{functionColor},
  moredelim =[is][\color{blue}]{-:}{:-},
  moredelim =[s][\color{accent}]{<}{>},
  moredelim=**[s][\color{black}]{/}{\ },
  moredelim=*[s][\color{blue}]{\$}{\ },
  alsoletter=:{:,-},
  literate = {(}{{\textcolor{black}{(}}}{1} {)}{{\textcolor{black}{)}}}{1} 
             {\{}{{\textcolor{black}{\{}}}{1} {\}}{{\textcolor{black}{\}}}}{1} 
             {last(}{{\textcolor{functionColor}{last}(}}5 {@<}{{<}}1
}


\lstdefinestyle{SQL}{
  style=cmput391,
  language=SQL,
  sensitive={True},
  deletekeywords={year,Cast,MIN,MAX},
  morekeywords={ABORT,AFTER,BEFORE,REFERENCES,REFERENCING,WITH,COMMITTED,UNCOMMITTED,REPEATABLE,SERIALIZABLE},
  keywords=[2]{AVG,MIN,MAX,COUNT,SUM,NEW,OLD,strftime,date,length,RAISE,rowid,substr},
  keywords=[3]{BIGINT,CHAR,DATE,DECIMAL,INT,FLOAT,TEXT},
  keywordstyle=[1]\ttfamily\bfseries\color{sqlColor},
  keywordstyle=[2]\ttfamily\bfseries\color{functionColor},
  keywordstyle=[3]\ttfamily\bfseries\color{datatypeColor},
  moredelim=**[is][\ttfamily\bfseries\color{functionColor}]{(@}{@)},
  morestring=[s][\color{stringColor}]{'}{'},
  moredelim=[is][\bfseries\color{highlight}]{-:}{:-}
}


%
% lstlisting for command line examples
%
\lstdefinestyle{commandLine}{
  basicstyle=cmput391,
  keywordstyle=[1]\ttfamily\bfseries\color{functionColor},
  keywordstyle=[2]\ttfamily\bfseries\color{datatypeColor},
  moredelim=**[is][\color{accent}\bf]{@}{@}
  keywords=[1]{%
    grep},
  keywords=[2]{%
    postings}
}

%
% lstlisting for C code
%
\lstdefinestyle{C}{
  language=C,
  tabsize=2,
  basicstyle=\ttfamily\footnotesize,
  stringstyle=\color{stringColor},
  keywords=[2]{sqlite3_prepare_v2,sqlite3_bind_int64,sqlite3_step,sqlite3_exec},
  keywordstyle=[1]\ttfamily\bfseries\color{datatypeColor},
  keywordstyle=[2]\ttfamily\bfseries\color{functionColor},
  commentstyle=\ttfamily\itshape\color{commentColor},
  moredelim=**[is][\color{Cfunction}\bf\ttfamily]{@}{@}
}


\lstdefinestyle{Python}{
  language=Python,
  tabsize=2,
  basicstyle=\ttfamily\footnotesize,
  stringstyle=\color{stringColor},
  morekeywords=[1]{self,yield,Class,true,false,None},
  keywords=[2]{satisfies,from,Iterator,EOF,concatenate,advance,join,rewind},
  deletekeywords=[1]{from},
  keywordstyle=[1]\ttfamily\bfseries\color{datatypeColor},
  keywordstyle=[2]\ttfamily\bfseries\color{functionColor},
  commentstyle=\ttfamily\itshape\color{commentColor},
  moredelim=**[is][\color{Cfunction}\bf\ttfamily]{@}{@}
}



%
% lstlisting style for RDF/SPARQL
%
\lstdefinestyle{SPARQL}{
  style=cmput391,
  sensitive={True},
  keywords=[1]{BASE,PREFIX,SELECT,WHERE,FILTER,MINUS,NOT,EXISTS},
  keywordstyle=[1]\ttfamily\bfseries\color{sqlColor},
  moredelim=*[s][\color{Cfunction}]{?}{\ },
  moredelim=*[s][\color{Cfunction}]{@}{\ },
  morestring=[s][\color{datatypeColor}]{"}{"},
  literate =* {|}{{|}}1 {//}{{//}}2 {< }{< }{2}
}

%
% coloring of markup code
%
\lstdefinestyle{markup}{
  style=cmput391,
  literate =* {=}{{\color{sqlColor}{=}}}1,
  moredelim =* [s][\color{accent}]{<}{>},
  moredelim =** [is][\color{accent}]{@}{@},
  moredelim =** [is][\color{blue}]{-:}{:-},
  morecomment=[s]{<!--}{-->}
}

\lstdefinestyle{RDF}{
  style=cmput391,
  keywords=[1]{BASE,PREFIX},
  keywordstyle=[1]\ttfamily\bfseries\color{functionColor},
  alsoletter=:{:,-,@},
  literate = {prefix}{{\bfseries\textcolor{functionColor}{prefix}}}6 
             {:}{{\textcolor{functionColor}{:}}}1 
             {@}{{\textcolor{functionColor}{@}}}1,
  moredelim =* [s][\color{accent}]{<}{>},
  moredelim =** [is][\color{functionColor}]{-:}{:-},
  morecomment=[s]{<!--}{-->}
}




















%
% Tikz stuff, for cummulative diagrams and algebra trees
%
\usepackage{tikz} %commutative diagrams
\usepackage{tikz-cd} %commutative diagrams
\usetikzlibrary{positioning}% To get more advances positioning options
\usetikzlibrary{arrows,shapes,automata}
\usetikzlibrary{calc} % for calculations and hard-coded coordinates
\usetikzlibrary{tikzmark, fit, shapes.misc}
\usetikzlibrary{decorations.pathreplacing}
\usetikzlibrary{backgrounds} % to draw underneath other elements
\usepackage{stackengine} % to overlay tikz pictures over other things

\usepackage{pgfplots}
\pgfplotsset{compat=1.17}

%
%
% PGF arrowtip that looks like an X
%
%
\pgfarrowsdeclare{X}{X}
{
  \arrowsize=0.25pt
  \advance\arrowsize by .5\pgflinewidth
  \pgfarrowsleftextend{-4\arrowsize-.5\pgflinewidth}
  \pgfarrowsrightextend{.5\pgflinewidth}
}
{
  \arrowsize=0.2pt
  \advance\arrowsize by .5\pgflinewidth
  \pgfsetdash{}{0pt} % do not dash
  \pgfsetroundjoin   % fix join
  \pgfsetroundcap    % fix cap
  \pgfpathmoveto{\pgfpointorigin}
  \pgfpathlineto{\pgfpoint{-5\arrowsize}{5\arrowsize}}
  \pgfusepathqstroke
  \pgfpathmoveto{\pgfpointorigin}
  \pgfpathlineto{\pgfpoint{5\arrowsize}{-5\arrowsize}}
  \pgfusepathqstroke
  \pgfpathmoveto{\pgfpointorigin}
  \pgfpathlineto{\pgfpoint{-5\arrowsize}{-5\arrowsize}}
  \pgfusepathqstroke
  \pgfpathmoveto{\pgfpointorigin}
  \pgfpathlineto{\pgfpoint{5\arrowsize}{5\arrowsize}}
  \pgfusepathqstroke
}


%
%%% for algorithms and pseudocode
%
\usepackage{algorithm}
% \usepackage{algpseudocode}
\usepackage[noend]{algpseudocode}
\renewcommand{\algorithmicrequire}{\textbf{Input:}}
\renewcommand{\algorithmicensure}{\textbf{Output:}}
\algnewcommand\algorithmicforeach{\textbf{for each}}
\algdef{S}[FOR]{ForEach}[1]{\algorithmicforeach\ #1\ \algorithmicdo}
% style for comments inside algorithms
\renewcommand{\algorithmiccomment}[1]{\hspace*{1em}\textcolor{commentColor}{\%~\textit{#1}}} 

%
% to create boxes with verbatim and code-like content
%
\usepackage{verbatimbox}

%
% to crop/trim boxes!
%
\usepackage{trimclip}

%
% colored and framed boxes
%
\newenvironment{BOX}[1]{\begin{tcolorbox}[colback=snow!25,colframe=snow!40!black,title=#1]\small}{\end{tcolorbox}}



%
% math and symbols
%
\usepackage{mathtools}
\DeclarePairedDelimiter{\ceil}{\lceil}{\rceil}
\DeclarePairedDelimiter{\floor}{\lfloor}{\rfloor}
\usepackage{amsbsy} 
\usepackage{amssymb}
\usepackage{ifsym} % needed to produce the common outerjoin symbols

%
% for nice-looking enumerations and inlined enumerations
%
\usepackage[shortlabels,inline]{enumitem}



%%% outer join symbols
\newcommand\TextbookOuterJoin{%
  \mathrel{\ooalign{\hss$\Join$\hss\cr%
  \kern0.5ex\raise1ex\hbox{\scalebox{0.7}{$\circ$}}}}}

\newcommand\RightOuterJoin{%
  \mathrel{%
  	\ooalign{%
  		$\Join$\cr\kern1.25ex\raise0.0975ex\hbox{\scalebox{0.25}[0.8275]{$\sqsubset$}}}
  	}
}

\newcommand\LeftOuterJoin{%
  \mathrel{%
  	\ooalign{%
  		$\Join$\cr\kern-0.0125ex\raise0.0975ex\hbox{\scalebox{0.25}[0.8275]{$\sqsupset$}}}%
	}
}

\newcommand\FullOuterJoin{%
  \mathrel{%
  	\ooalign{%
  		$\Join$\cr\kern-0.0125ex\raise0.0975ex\hbox{\scalebox{0.25}[0.8275]{$\sqsupset$}}}%
  		\kern-0.45ex\raise0.0975ex\hbox{\scalebox{0.25}[0.8275]{$\sqsubset$}}
	}
}




